% ============ 2 平展态射 (原书 pp.16-24) ============
\chapter{平展态射}\label{sec:2}
\index{平展态射}
\index{切空间}
\index{切锥}
\index{平坦态射}
\index{非分歧态射}
\index{标准平展态射}
\index{拟有限态射}
\index{局部等平凡}

平展态射是代数几何中如下对象的类比：微分几何中流形的局部同构，复分析中没有分支点的 Riemann 曲面的覆盖，代数数论中的非分歧扩张.

对簇，可以几何地刻画平展态射；对任意概形，只有交换代数.

\subsection*{非奇异代数簇的平展态射}

本节中，所有簇都定义在代数闭域 $k$ 上.设 $W$ 与 $V$ 是 $k$ 上的非奇异代数簇.如引言中那样，正则映射 $\varphi \colon W \to V$ 称为在 $Q \in W$ 处平展，如果切空间上的映射 $d\varphi \colon \operatorname{Tgt}_{Q}(W) \to \operatorname{Tgt}_{\varphi(Q)}(V)$ 是同构，而 $\varphi$ 称为平展的，如果它在 $W$ 的每一点处都平展.

\begin{proposition}\label{prop:2.1}
设 $V = \operatorname{Specm}(A)$ 是 $k$ 上的非奇异仿射簇，$W$ 是由方程
\[
g_{i}(Y_1, \ldots, Y_n) = 0, \qquad g_{i} \in A[Y_1, \ldots, Y_n], \quad i = 1, \ldots, n
\]
定义的 $V \times \mathbb{A}^{n}$ 的子簇.投影映射 $W \to V$ 在 $W$ 的点 $(P; b_1, \ldots, b_n)$ 处平展，当且仅当 Jacobi 矩阵 $\left(\frac{\partial g_i}{\partial Y_j}\right)$ 在 $(P; b_1, \ldots, b_n)$ 处非奇异.
\end{proposition}

\begin{proof}
置
\[
A = k[X_1, \ldots, X_m]/(F_1, \ldots, F_r)
\]
且
\[
P = (a_1, \ldots, a_m).
\]
$V$ 在 $P$ 处的切空间是线性方程组
\[
(dF_i)_{P} \overset{\mathrm{def}}{=} \sum_{j=1}^{m} \left.\frac{\partial F_i}{\partial X_j}\right|_{P} (dX_j)_{P} = 0, \qquad i = 1, \ldots, r
\]
在变量 $(dX_j)_{P}$ 中的解空间(见 AG \S 5).

按定义，$W = \operatorname{Specm} A[Y_1, \ldots, Y_n]/(g_1, \ldots, g_n)$.对每个 $i$，选取 $G_i \in k[X_1, \ldots, X_m, Y_1, \ldots, Y_n]$ 映到 $g_i \in A[Y_1, \ldots, Y_n]$.则
\[
A[Y_1, \ldots, Y_n]/(g_1, \ldots, g_n) = k[X_1, \ldots, X_m, Y_1, \ldots, Y_n]/(F_1, \ldots, F_r, G_1, \ldots, G_n),
\]
从而 $W$ 在 $(P; b_1, \ldots, b_n)$ 处的切空间是线性方程组
\begin{align*}
\sum_{j=1}^{m} \left.\frac{\partial F_i}{\partial X_j}\right|_{P} (dX_j)_{(P;b)} &= 0, \quad i = 1, \ldots, r; \\
\sum_{j=1}^{m} \left.\frac{\partial G_i}{\partial X_j}\right|_{(P;b)} (dX_j)_{P} + \sum_{j=1}^{n} \left.\frac{\partial G_i}{\partial Y_j}\right|_{(P;b)} (dY_j)_{(P;b)} &= 0, \quad i = 1, \ldots, n
\end{align*}
($m + n$ 个变量中的 $r + n$ 个方程)的解空间.映射 $\operatorname{Tgt}_{(P;b)}(W) \to \operatorname{Tgt}_{P}(V)$ 是显然的投影映射.因此，$\varphi$ 在 $(P; b)$ 处平展，当且仅当第一个方程组的每个解都唯一地延拓为第二个方程组的解.当且仅当 $n \times n$ 矩阵 $\left.\left(\frac{\partial G_i}{\partial Y_j}\right)\right|_{(P;b)}$ 非奇异时，才会如此.
\end{proof}

\begin{corollary}\label{cor:2.2}
设 $\varphi \colon U \to V$ 是正则映射，其中 $U$ 与 $V$ 都等于 $\mathbb{A}^{m}$.则 $\varphi$ 在 $(a_1, \ldots, a_m)$ 处平展，当且仅当 Jacobi 矩阵 $\left.\left(\frac{\partial (X_i \circ \varphi)}{\partial Y_j}\right)\right|_{(a_1, \ldots, a_m)}$ 非奇异.(这里 $X_i$ 是 $V$ 上的第 $i$ 个坐标函数，$Y_j$ 是 $U$ 上的第 $j$ 个坐标函数.)
\end{corollary}

\begin{proof}
每个 $X_i \circ \varphi$ 都是 $U = \mathbb{A}^{m}$ 上的正则函数，从而是 $U$ 上坐标函数 $Y_j$ 的多项式 $G_i(Y_1, \ldots, Y_m)$.$\varphi$ 的图象 $\Gamma_{\varphi}$ 是由方程
\[
X_i = G_i(Y_1, \ldots, Y_m), \qquad i = 1, \ldots, m
\]
定义的 $\mathbb{A}^{m} \times \mathbb{A}^{m}$ 的子簇.映射 $\varphi$ 是同构 $P \mapsto (P, \varphi(P)) \colon V \to \Gamma_{\varphi}$(见 AG 4.26)与投影 $\Gamma_{\varphi} \to \mathbb{A}^{m}$ 的复合，从而在 $P_0 = (a_1, \ldots, a_m)$ 处平展，当且仅当投影映射在 $(P_0, \varphi(P_0))$ 处平展.命题现在给出论断.
\end{proof}

\begin{remark}\label{rem:2.3}
注意，条件``Jacobi 矩阵 $\left.\left(\frac{\partial (X_i \circ \varphi)}{\partial Y_j}\right)\right|_{(a_1, \ldots, a_m)}$ 非奇异''正是高等微积分中逆映射定理的假设.
\end{remark}

\begin{exercise}\label{exc:2.4}
(a) 用对偶数意义下的切空间定义证明该命题(见 AG 5.37).

(b) 把 $U$ 换成 $\mathbb{A}^{n}$ 的开子集证明该推论.
\end{exercise}

\begin{example}\label{ex:2.5}
设 $V = \operatorname{Specm}(A)$ 是 $k$ 上的非奇异仿射簇，
\[
f(T) = a_0 T^{m} + \cdots + a_m
\]
是系数在 $A$ 中的多项式.于是每个 $a_i$ 都是 $V$ 上的正则函数，我们可以把 $f(T)$ 看作由 $P \in V$ 参数化的多项式族 $f(P; T)$.设 $W = \operatorname{Specm}(A[T]/(f(T)))$(假定该环是既约的).则 $W$ 是由方程
\[
f(P; T) = 0
\]
定义的 $V \times \mathbb{A}^{1}$ 的子簇，而包含 $A \hookrightarrow A[T]/(f(T))$ 对应于投影映射 $\pi \colon W \to V$, $(P, c) \mapsto P$.对每个 $P_0 \in V$，$\pi^{-1}(P_0)$ 是 $f(P_0; T)$ 的根集，其中
\[
f(P_0; T) = a_0(P_0) T^{m} + \cdots + a_m(P_0) \in k[T].
\]
我们有：
\begin{enumerate}
  \item 纤维 $\pi^{-1}(P_0)$ 有限，当且仅当 $P_0$ 不是诸 $a_i$ 的公共零点；从而 $\pi$ 拟有限(AG p.8.5)，当且仅当由 $a_0, a_1, \ldots, a_m$ 生成的理想是 $A$；
  \item 映射 $\pi$ 有限，当且仅当 $a_0$ 是 $A$ 中的单位(比较 AG 8.1).
  \item 映射 $\pi$ 在 $(P_0; c)$ 处平展，当且仅当 $c$ 是 $f(P_0; T)$ 的单根.
\end{enumerate}
为证明 (c)，注意 $c$ 是 $f(P_0; T)$ 的单根，当且仅当 $c$ 不是 $\frac{df(P_0; T)}{dT}$ 的根.现在可以应用命题~\ref{prop:2.1}.
\end{example}

\begin{example}\label{ex:2.6}
考虑映射 $x \mapsto x^{n} \colon \mathbb{A}^{1} \to \mathbb{A}^{1}$.由于 $\frac{dX^{n}}{dX} = nX^{n-1}$，由推论~\ref{cor:2.2} 看到，若 $k$ 的特征整除 $n$，则该映射在 $\mathbb{A}^{1}$ 的任何一点都不平展，否则它在所有 $x \neq 0$ 处平展.
\end{example}

\subsection*{任意簇的平展态射}

簇的奇异点处的切空间对了解该点几乎无用.相反，必须使用切锥(AG \S 5).

回忆对代数闭域 $k$ 上的仿射簇 $V = \operatorname{Specm} k[X_1, \ldots, X_n]/\mathfrak{a}$，原点 $P$ 处的切锥由理想 $\mathfrak{a}^{*} = \{f^{*} \mid f \in \mathfrak{a}\}$ 定义，其中对 $f \in k[X_1, \ldots, X_n]$，$f^{*}$ 是 $f$ 的最低次齐次部分.$P$ 处的几何切锥是 $\mathfrak{a}^{*}$ 的零集.然而，$k[X_1, \ldots, X_n]/\mathfrak{a}^{*}$ 可能有幂零元，因此(如下面的例~\ref{ex:2.7}(b))几何切锥可能不决定这个环.作为替代，我们把切锥 $C_{P}(V)$ 定义为 $k$-代数 $k[X_1, \ldots, X_n]/\mathfrak{a}^{*}$(或者等价地但更好，仿射 $k$-概形 $\operatorname{Spec} k[X_1, \ldots, X_n]/\mathfrak{a}^{*}$).

设 $\varphi \colon W \to V$ 是代数闭域 $k$ 上的簇的正则映射.则 $\varphi$ 称为在 $Q \in W$ 处平展，如果它诱导切锥的同构 $C_{\varphi(Q)}(V) \to C_{Q}(W)$(作为 $k$-代数).对非奇异簇，这与上一小节的定义一致.

切锥可以更典范地定义.回忆(Atiyah 和 MacDonald 1969, p.111)每个局部环 $A$ 都联系一个分次环
\[
\operatorname{gr}(A) \overset{\mathrm{def}}{=} \bigoplus_{n} \mathfrak{m}^{n}/\mathfrak{m}^{n+1}, \qquad \mathfrak{m} = \mathfrak{m}_{A}.
\]
$\operatorname{gr}(A)$ 上的乘法由乘法
\[
a, b \mapsto ab \colon \mathfrak{m}^{i} \times \mathfrak{m}^{j} \to \mathfrak{m}^{i+j}
\]
诱导.

局部同态 $A \to B$ 定义分次环的同态 $\operatorname{gr}(A) \to \operatorname{gr}(B)$.容易证明，若 $\operatorname{gr}(A) \to \operatorname{gr}(B)$ 是同构，则完备化上诱导的映射 $\widehat{A} \to \widehat{B}$ 也是同构(同上，10.23).逆命题也真，因为 $\operatorname{gr}(A) = \operatorname{gr}(\widehat{A})$(同上，10.22).\footnote{回忆我们总假定所有环都是 Noether 的.}

\begin{example}\label{ex:2.7}
(a) 曲线
\[
V \colon Y^{2} = X^{3} + X^{2}
\]
在原点 $P$ 处的切锥由方程
\[
Y^{2} = X^{2}
\]
定义.因此它是直线 $Y = \pm X$ 的并.更精确地说，设 $x$ 与 $y$ 分别表示 $X$ 与 $Y$ 在 $\mathfrak{m}_{P}/\mathfrak{m}^{2}_{P}$ 中的类；则 $C_{P}(V) = k[x, y]/(y^{2} - x^{2})$.考虑映射 $t \mapsto (t^{2} - 1, t(t^{2} - 1)) \colon \mathbb{A}^{1} \to V$.它对应于 $k$-代数映射
\[
X \mapsto T^{2} - 1, \quad Y \mapsto T(T^{2} - 1) \colon k[X, Y]/(Y^{2} - X^{2} - X^{3}) \to k[T].
\]
$\mathbb{A}^{1}$ 在 $Q = 1$ 处的切锥是 $k[s]$，其中 $s$ 表示 $T - 1$ 在 $\mathfrak{m}_{Q}/\mathfrak{m}^{2}_{Q}$ 中的类.映射 $C_{P} \to C_{Q}$ 是 $k$-代数映射
\[
x \mapsto 2s, \quad y \mapsto 2s \colon k[x, y]/(y^{2} - x^{2}) \to k[s].
\]
这不是同构，因此该映射在 $Q$ 处不平展.

(b) 曲线
\[
V \colon Y^{2} = X^{3}
\]
在原点处的切锥由方程
\[
Y^{2} = 0
\]
定义.因此它是带重数 2 的直线 $Y = 0$.映射 $t \mapsto (t^{2}, t^{3}) \colon \mathbb{A}^{1} \to V$ 在原点不平展，因为它在切锥上定义的映射
\[
x \mapsto 0, \quad y \mapsto 0 \colon k[x, y]/(y^{2}) \to k[t]
\]
不是同构.
\end{example}

对代数闭域 $k$ 上的簇 $V$ 上的任意点 $P$，$C_{P}(V) = \operatorname{gr}(\mathcal{O}_{P})$(AG 5.40).于是我们得到如下判别准则：

\begin{quote}
代数闭域 $k$ 上的簇的正则映射 $\varphi \colon W \to V$ 在 $W$ 的点 $Q$ 处平展，当且仅当 $\varphi$ 诱导的映射 $\widehat{\mathcal{O}}_{V, \varphi(Q)} \to \widehat{\mathcal{O}}_{W, Q}$ 是同构.
\end{quote}

特别地，若 $Q$ 映到 $P$，则 $\varphi$ 在 $Q$ 处平展的一个必要条件是：$V$ 在 $P$ 处必须有与 $W$ 在 $Q$ 处相同类型的奇异性.

\textbf{任意域上的簇的平展态射} 设 $\varphi \colon W \to V$ 是域 $k$ 上的簇的正则映射.我们说 $\varphi$ 在 $w \in W$ 处平展，如果对 $k$ 的某个代数闭包 $k^{\mathrm{al}}$，$\varphi_{k^{\mathrm{al}}} \colon W_{k^{\mathrm{al}}} \to V_{k^{\mathrm{al}}}$ 在 $W_{k^{\mathrm{al}}}$ 中映到 $w$ 的那些点处平展.

\begin{exercise}\label{exc:2.8}
用 $W$ 与 $V$ 的切空间(非奇异情形)或切锥(一般情形)使这个条件显式化.证明(或否证)：$\varphi \colon W \to V$ 在 $w$ 处平展，当且仅当映射 $\operatorname{gr}(\mathcal{O}_{V, v}) \otimes_{\kappa(v)} \kappa(w) \to \operatorname{gr}(\mathcal{O}_{W, w})$($v = \varphi w$)是同构.
\end{exercise}

\subsection*{概形的平展态射}

本小节中论断的证明见 EC I.1--I.3.

\textbf{平坦态射} 回忆环同态 $A \to B$ 是平坦的，如果函子 $M \mapsto B \otimes_{A} M$(从 $A$-模到 $B$-模)正合.也称 $B$ 是平坦 $A$-代数.要检验 $f \colon A \to B$ 是否平坦，只需检验对 $B$ 的每个极大理想 $\mathfrak{m}$，局部同态 $A_{f^{-1}(\mathfrak{m})} \to B_{\mathfrak{m}}$ 平坦.

若 $A$ 是整环，则对一切非零 $a$，$x \mapsto ax \colon A \to A$ 是单射.因此，对任意平坦 $A$-代数 $B$，$x \mapsto ax \colon B \to B$ 也是单射，从而 $A \to B$ 是单射.对 Dedekind 整环 $A$，逆命题成立：同态 $A \to B$ 平坦当且仅当它单射.

概形(或簇)的态射 $\varphi \colon Y \to X$ 是平坦的，如果对所有 $y \in Y$，局部同态 $\mathcal{O}_{X, \varphi(y)} \to \mathcal{O}_{Y, y}$ 平坦.平坦性定义后的评注表明，只需对闭点 $y \in Y$ 检验这一点.

簇的平坦态射 $\varphi \colon Y \to X$ 是代数几何中如下对象的类比：由 $X$ 的点参数化的流形族 $Y_x \overset{\mathrm{def}}{=} \varphi^{-1}(x)$ 的连续族.若 $\varphi$ 平坦，则
\[
\dim Y_x = \dim Y - \dim X
\]
对所有使 $Y_x \neq \varnothing$ 的(闭)点 $x \in X$ 成立；若 $X$ 与 $Y$ 非奇异，逆命题也真.此外，簇的有限态射 $\varphi \colon Y \to X$ 是平坦的，当且仅当每个纤维 $\varphi^{-1}(x)$ 按重数计都有相同数目的点.(若 $\varphi$ 是由仿射 $k$-代数同态 $A \to B$ 定义的仿射簇映射，则条件是说 $\dim_{A/\mathfrak{m}} B/\mathfrak{m}B$ 与极大理想 $\mathfrak{m} \subset A$ 无关.)

设 $Z$ 是 $X$ 的闭子概形.则包含 $Z \hookrightarrow X$ 平坦，当且仅当 $Z$ 在 $X$ 中也是开的(从而是 $X$ 的连通分支).

\textbf{非分歧态射} 局部环的局部同态 $f \colon A \to B$ 是非分歧的，如果 $B/f(\mathfrak{m}_{A})B$ 是 $A/\mathfrak{m}_{A}$ 的有限可分域扩张，或者等价地，如果
\begin{enumerate}
  \item $f(\mathfrak{m}_{A})B = \mathfrak{m}_{B}$，且
  \item 域 $B/\mathfrak{m}_{B}$ 在 $A/\mathfrak{m}_{A}$ 上有限且可分.
\end{enumerate}
这与代数数论中的定义一致，在那里只考虑离散赋值环.

概形的态射 $\varphi \colon Y \to X$ 是非分歧的，如果它有限型，并且对所有 $y \in Y$，映射 $\mathcal{O}_{X, f(y)} \to \mathcal{O}_{Y, y}$ 非分歧.只需对 $Y$ 的闭点 $y$ 检验该条件.

设 $\varphi \colon Y \to X$ 是有限型态射.则 $\varphi$ 非分歧，当且仅当层 $\Omega^{1}_{Y/X}$ 为零(见 EC I.3.5).

\textbf{平展态射} 概形的态射 $\varphi \colon Y \to X$ 是平展的，如果它平坦且非分歧.特别地，这意味着 $\varphi$ 有限型.\footnote{Grothendieck 只要求它局部有限型，或者对任意(不一定是 Noether 的)概形要求它局部有限表示.}

环同态 $f \colon A \to B$ 是平展的，如果 $\operatorname{Spec} B \to \operatorname{Spec} A$ 平展.等价地，它平展，如果
\begin{enumerate}
  \item $B$ 是有限生成的 $A$-代数；
  \item $B$ 是平坦 $A$-代数；
  \item 对 $B$ 的所有极大理想 $\mathfrak{n}$，$B_{\mathfrak{n}}/f(\mathfrak{p})B_{\mathfrak{n}}$ 是 $A_{\mathfrak{p}}/\mathfrak{p}A_{\mathfrak{p}}$ 的有限可分域扩张，其中 $\mathfrak{p} = f^{-1}(\mathfrak{n})$.
\end{enumerate}

设 $X = \operatorname{Spec} A$，其中 $A$ 是整环.对任意真理想 $\mathfrak{a} \subset A$，对应于同态 $A \to A/\mathfrak{a}$ 的映射 $Z \hookrightarrow X$ 非分歧，但不平坦，从而不平展.这与``平展''意为``局部同构''的直觉一致：把真闭子流形包含进连通流形不是局部同构.

\begin{proposition}\label{prop:2.9}
对代数闭域上簇的正则映射 $\varphi \colon Y \to X$，本小节中``平展''的定义与上一小节中的一致.
\end{proposition}

\begin{proof}
首先注意，簇的每个正则映射都有限型.因此，只需证明：对由 $k$-簇的正则映射产生的每个局部 $k$-代数局部同态 $f \colon A \to B$，同态 $\widehat{f} \colon \widehat{A} \to \widehat{B}$ 是同构，当且仅当 $f$ 平坦且非分歧.

当然，若 $\widehat{A} \to \widehat{B}$ 是同构，则它既平坦又非分歧.因此为证必要性，只需证明：若 $\widehat{A} \to \widehat{B}$ 平坦或非分歧，则 $A \to B$ 有同样的性质.若 $\widehat{A} \to \widehat{B}$ 平坦，则 $A$-模的正合序列 $M' \to M \to M''$ 给出 $\widehat{B}$-模的正合序列
\[
\widehat{B} \otimes_{A} M' \to \widehat{B} \otimes_{A} M \to \widehat{B} \otimes_{A} M''
\]
(因为 $\widehat{A}$ 是平坦 $A$-代数---Atiyah 和 MacDonald 1969, 10.14)，这蕴含
\[
B \otimes_{A} M' \to B \otimes_{A} M \to B \otimes_{A} M''
\]
正合(因为 $\widehat{B}$ 是忠实平坦 $B$-代数，同上，10，习题 7).若 $\widehat{A} \to \widehat{B}$ 非分歧，则理想 $\mathfrak{m}_{B}$ 与 $\mathfrak{m}_{A}B$ 都生成 $\widehat{B}$ 中的极大理想，这蕴含它们相等(因为对环 $R$ 中的任意理想 $\mathfrak{a}$ 与忠实平坦同态 $R \to S$，有 $\mathfrak{a}S \cap R = \mathfrak{a}$\footnote{Matsumura, H., \emph{Commutative Algebra}, Benjamin, 1970, 4.C.}).

逆命题不那么初等.Zariski 主定理(EC I 1.8)使我们能够假定 $B$ 是有限 $A$-代数 $C$ 在某个极大理想处的局部化，并且事实上 $C = A[T]/(f(T))$，其中 $f(T)$ 是某个首一多项式.现在 $\widehat{B}$ 是 $\widehat{A}[T]/(f(T))$ 在位于 $\mathfrak{m}\widehat{A}$ 之上的某个极大理想处的完备化.由于 Hensel 引理(Atiyah 和 MacDonald 1969, 10，习题 9)，在 $k[T]$ 中($k = A/\mathfrak{m}_{A}$)的分解 $\overline{f} = g^{e_1}_{1} \cdots g^{e_n}_{n}$，其中 $g_i$ 是 $\overline{f}$ 的互异不可约因子，提升为 $\widehat{A}[T]$ 中的分解 $f = f_1 \cdots f_n$.相应地，$\widehat{A}[T]/(f(T)) = \prod_{i} \widehat{A}[T]/(f_i(T))$，从而对某个 $i$，$\widehat{B} = \widehat{A}[T]/(f_i(T))$.于是 $\widehat{A} \to \widehat{B}$ 有限，因此可以应用下一个引理.
\end{proof}

\textbf{注记} 平坦非分歧局部环同态 $A \to B$($A, B$ 的剩余域相同)在完备化上诱导同构的另一证明.由平坦性的一个结果(Bourbaki, \emph{Commutative Algebra}, p.227; Matsumura, \emph{Commutative Ring Theory}, p.74)，诱导映射 $\mathfrak{m}^{n}_{A}/\mathfrak{m}^{n+1}_{A} \to \mathfrak{m}^{n}_{B}/\mathfrak{m}^{n+1}_{B}$ 是同构.因此映射 $A/\mathfrak{m}^{n}_{A} \to B/\mathfrak{m}^{n}_{B}$ 是同构.对 $n$ 取极限，我们发现完备化上的映射是同构.

\begin{lemma}\label{lem:2.10}
设 $\varphi \colon A \to B$ 是局部环的局部同态.若
\begin{enumerate}
  \item $\varphi$ 单射，
  \item 剩余域上的映射 $A/\mathfrak{m}_{A} \to B/\mathfrak{m}_{B}$ 是同构，
  \item $\varphi$ 非分歧，且
  \item $B$ 是有限 $A$-代数，
\end{enumerate}
则 $\varphi$ 是同构.
\end{lemma}

\begin{proof}
由 (a)，只需证明 $\varphi$ 满射.由 (b) 和 (c) 看到
\[
B = \varphi(A) + \mathfrak{m}_{B} = \varphi(A) + \mathfrak{m}_{A}B.
\]
条件 (d) 允许我们把 Nakayama 引理应用于 $A$-模 $B$，得到 $B = \varphi(A)$.
\end{proof}

\subsection*{例}

\textbf{Jacobi 判别法} 命题~\ref{prop:2.1} 对仿射概形 $V$ 成立.

\textbf{域} 设 $k$ 是域.局部 $k$-代数 $A$ 非分歧，当且仅当它是 $k$ 的有限可分域扩张.设 $A$ 是平展 $k$-代数.由于 $A$ 在 $k$ 上非分歧，$A$ 的极大理想都不含真素理想，因此 $A$ 的维数为 0.从而它是 Artin 环(Atiyah 和 MacDonald 1969, 8.5)，因此是 $k$ 的有限可分域扩张的有限乘积(同上，8.7).反过来，$k$ 的有限可分域扩张的有限乘积是平展 $k$-代数.

\textbf{Dedekind 整环} 设 $A$ 是 Dedekind 整环，$L$ 是其分式域 $K$ 的有限可分扩张.设 $B$ 是 $A$ 在 $L$ 中的整闭包，$P$ 是 $B$ 的素理想.则 $\mathfrak{p} \overset{\mathrm{def}}{=} P \cap A$ 是 $A$ 的素理想，而 $P$ 称为在 $B/A$ 中非分歧，如果
\begin{enumerate}
  \item 在 $\mathfrak{p}B$ 分解为素理想乘积的过程中，$P$ 出现的指数为 1；
  \item 剩余域扩张 $B/P \supset A/\mathfrak{p}$ 可分.
\end{enumerate}
这等价于映射 $A_{\mathfrak{p}} \to B_{P}$ 非分歧.设 $b$ 是 $B$ 中属于 $B$ 的所有分歧素理想的元素.则 $B_b \overset{\mathrm{def}}{=} B[b^{-1}]$ 是平展 $A$-代数，且每个平展 $A$-代数都是这种类型的代数的有限乘积.

\textbf{标准平展态射} 设 $A$ 是环，$f(T)$ 是系数在 $A$ 中的首一多项式.由于 $f(T)$ 首一，$A[T]/(f(T))$ 是有限秩的自由 $A$-模，从而是平坦的.对任意使 $f'(T)$ 在 $(A[T]/(f(T)))_{b}$ 中可逆的 $b \in A[T]/(f(T))$，同态 $A \to (A[T]/(f(T)))_{b}$ 平展.平展态射 $\varphi \colon V \to U$ 称为标准的，如果它同构于这样一个同态的谱，即如果存在环 $A$，首一多项式 $f(T) \in A[T]$ 和满足上述条件的 $b \in A[T]/(f(T))$，使得有交换图：
\[
\begin{tikzcd}[row sep=small]
  V \arrow[r] \arrow[rd] & \operatorname{Spec}(A[T]/(f(T)))_{b} \arrow[d] \\
                         & \operatorname{Spec} A
\end{tikzcd}
\]
其中 $\varphi$ 与 $\pi$ 如所示.一个重要事实是，局部地每个平展态射都是标准的，即对任意平展态射 $\varphi \colon Y \to X$ 与 $y \in Y$，存在 $y$ 与 $\varphi(y)$ 的开仿射邻域 $V$ 与 $U$，使得 $\varphi(V) \subset U$ 且 $\varphi|_{V} \colon V \to U$ 是标准的(EC I 3.14).

\textbf{正规概形} 设 $X$ 是连通正规概形(或簇)；若 $X$ 仿射，则意味着 $X$ 是整闭整环 $A$ 的谱.设 $K$ 是 $X$ 上的有理函数域，$L$ 是 $K$ 的有限可分域扩张.设 $Y$ 是 $X$ 在 $L$ 中的正规化：若 $X = \operatorname{Spec} A$，则 $Y = \operatorname{Spec} B$，其中 $B$ 是 $A$ 在 $L$ 中的整闭包.则 $Y \to X$ 有限(Atiyah 和 MacDonald 1969, 5.17)，从而是有限型的.设 $U$ 是 $Y$ 的开子集，除去 $\varphi$ 分歧的闭点 $y$.则 $U \to X$ 平坦，从而是平展的；此外，每个平展 $X$-概形都是这种类型的平展 $X$-概形的不交并(EC I 3.20, 3.21).

\textbf{带幂零元的概形} 设 $X$ 是概形，$X_0$ 是由幂零理想定义的 $X$ 的子概形(于是 $X_0$ 与 $X$ 有相同的底层拓扑空间)；则映射 $U \mapsto U_0 \overset{\mathrm{def}}{=} U \times_{X} X_0$ 是平展 $X$-概形范畴到平展 $X_0$-概形范畴的等价.

\subsection*{平展态射的性质}

粗略地说，平展态射具有与局部同构的类比所提示的一切性质.这里我们只列出最重要的性质，并简要讨论其证明.

\begin{proposition}\label{prop:2.11}
(a) 每个开浸入都是平展的.

(b) 两个平展态射的复合是平展的.

(c) 平展态射的每个基变换都是平展的.

(d) 若 $\varphi \circ \psi$ 与 $\varphi$ 平展，则 $\psi$ 也平展.
\end{proposition}

论断 (a) 说，若 $U$ 是 $X$ 的开子概形(或子簇)，则包含 $U \hookrightarrow X$ 平展---这在任何定义下都是显然的.

论断 (b) 在任何定义下也是显然的.

关于簇的纤维积与正则映射的基变换的概念，见 (AG \S 4).给定点 $(u, y) \in U \times_{X} Y$，容易证明(例如用对偶数，同上，5.37)
\[
\operatorname{Tgt}_{(u,y)}(U \times_{X} Y) = \operatorname{Tgt}_{u}(U) \times_{\operatorname{Tgt}_{x}(X)} \operatorname{Tgt}_{y}(Y),
\]
其中 $x$ 是 $u$ 与 $y$ 在 $X$ 中的公共像.因此，对非奇异簇，(c) 是显然的.

最后，对簇，用切锥的定义证明 (d) 是显然的，例如.对概形，见 (EC I 3.6).

\begin{proposition}\label{prop:2.12}
设 $\varphi \colon Y \to X$ 是平展态射.

(a) 对所有 $y \in Y$，$\mathcal{O}_{Y, y}$ 与 $\mathcal{O}_{X, x}$ 有相同的 Krull 维数.

(b) 态射 $\varphi$ 拟有限.

(c) 态射 $\varphi$ 是开的.

(d) 若 $X$ 既约，则 $Y$ 也既约.

(e) 若 $X$ 正规，则 $Y$ 也正规.

(f) 若 $X$ 正则，则 $Y$ 也正则.
\end{proposition}

当 $X$ 与 $Y$ 是连通簇时，(a) 说它们有相同的维数---这是显然的，因为切锥与簇有相同的维数.

论断 (b) 说 $\varphi$ 的纤维都有限.

论断 (c) 由更一般的事实推出：有限型平坦态射是开的(EC I 2.12).

论断 (e) 由我们对到正规概形的平展态射的描述蕴含.

对簇，(f) 蕴含：若 $X$ 非奇异且 $Y \to X$ 平展，则 $Y$ 也非奇异.这是显然的，因为点 $P$ 非奇异，当且仅当 $P$ 处的切锥是 $\dim X$ 个变量的多项式环.

\begin{exercise}\label{exc:2.13}
找到非奇异簇的正则映射 $Y \to X$ 是开的直接简单证明.(这样的映射平坦，从而开；问题是只使用例如 AG 中证明的事实找到直接证明.)
\end{exercise}

\begin{proposition}\label{prop:2.14}
设 $\varphi \colon Y \to X$ 是有限型态射.$\varphi$ 平展的点集 $U \subset Y$ 在 $Y$ 中开.
\end{proposition}

当然，$U$ 可能是空的.若 $X$ 正规，则 $U$ 是 $\mathcal{O}_Y$-层 $\Omega^{1}_{Y/X}$ 的支集的补集.

\begin{proposition}\label{prop:2.15}
设 $\varphi \colon Y \to X$ 是簇的平展态射.若 $X$ 连通，则 $\varphi$ 的每个截面都是 $X$ 到 $Y$ 的某个连通分支上的同构.
\end{proposition}

\begin{proof}
为简单起见，假定基域代数闭.设 $s \colon X \to Y$ 是 $\varphi$ 的截面，即满足 $\varphi \circ s = \mathrm{id}_{X}$ 的正则映射.$\varphi$ 的图象 $\Gamma_{s} \overset{\mathrm{def}}{=} \{(x, s(x)) \mid x \in X\}$ 在 $X \times Y$ 中闭(见 AG 4.26)，它在正则映射 $y \mapsto (\varphi(y), y) \colon Y \to X \times Y$ 下的原像就是 $s(X)$.因此，$s(X)$ 在 $Y$ 中闭.它也是开的，因为 $s$ 平展(命题~\ref{prop:2.12}(c))，所以是 $Y$ 的连通分支.于是 $\varphi|_{s(X)}$ 与 $s$ 是互逆的同构.
\end{proof}

\begin{corollary}\label{cor:2.16}
设 $p \colon X \to S$ 与 $q \colon Y \to S$ 是代数闭域上簇的态射.假定 $p$ 平展且 $Y$ 连通.设 $\varphi, \varphi'$ 是满足 $p \circ \varphi = q$ 与 $p \circ \varphi' = q$ 的态射 $Y \to X$.若 $\varphi$ 与 $\varphi'$ 在 $Y$ 的某一点处一致，则它们在 $Y$ 上处处相等.
\end{corollary}

\begin{proof}
映射 $(1, \varphi), (1, \varphi') \colon Y \to Y \times_{S} X$ 是投影映射 $(y, x) \mapsto y \colon Y \times_{S} X \to Y$ 的截面，而假设说它们在某个点 $y_0 \in Y$ 处一致.由于投影映射平展，$(1, \varphi)$ 与 $(1, \varphi')$ 相等.但 $\varphi$ 与 $\varphi'$ 是 $(1, \varphi)$ 与 $(1, \varphi')$ 同投影映射 $X \times_{S} Y \to X$ 的复合，因此它们相等.
\end{proof}

\begin{remark}\label{rem:2.17}
若假定 $\varphi$ 分离，则命题~\ref{prop:2.15} 对概形成立.只要不仅假定 $\varphi$ 与 $\varphi'$ 在一点处一致，还假定它们在这一点诱导剩余域上的相同映射，推论对概形也成立.见 (EC I 3.12, I 3.13).
\end{remark}
